\documentclass{subfiles}
\begin{document}
    Evaluating polynomial in a projective space is tricky.
        Take for instance the following polynomial:
        \[
            f(X) = f(x_{0}, x_{1}, x_{2}) = x_{0} - x_{2}x_{1}.
        \]
        If we consider the points \(P = [2, -1, 1]\) and \(Q = 3P = [6, -3, 3]\),
            which are equal under the homogeneous interpretation,
            we get 
            \[
                f(P) = 2 - (-1 * 1) = 3 \text{ and } f(Q) = 6 - (-3 * 3) = 15.
            \]

    Polynomials like \(f\) are said to be non-homogeneous. 
        If instead we consider an homogeneous polynomial, we get 
        \[
            f(\lambda x_{0}, \lambda x_{1}, \lambda x_{2}) = \lambda^{d} f(x_{0}, x_{1}, x_{2})
        \]
        where \(d\) is the degree of the monomyals.
        It follows that, if \(f(X) \text{ and } g(X)\) are two homogeneous polynomials,
        then 
        \[
            F(X) = \frac{f(X)}{g(X)}
        \]
        does not depend by any proportionality factor.

    Luckly for us, in the projective space, 
        any polynomial can always be seen as the ratio of two homogeneous polynomials.
        For instance, the following non-homogeneous polynomial
        \[
            F(x, y) = 3x^{2} -2y + 1 
        \]
        we can transform it into an homogeneous polynomial as follow:
        we replace \(x\) by \(\sfrac{x_{1}}{x_{0}}\) and \(y\) by \(\sfrac{x_{2}}{x_{0}}\),
        compute the lowest common multiple and obtain
        \[
            3x^{2} - 2y + 1 \frac{3x_{1}^{2} - 2x_{0}x_{2} + x_{0}^{2}}{x_{0}^{2}},
        \]
        which evaluated at both \(P\) and \(Q\) give the same value of \(\sfrac{3}{4}\).
\end{document}
